\section*{Chapitre \ref{ch:g10:light-spectra} --- Spectres de la lumière et le message de la lumière}

\begin{solution}{exo:g10:light-spectra:1}
$\qty{254}{nm} < \qty{400}{nm}$~: \omterm{def:g10:light-spectra:wavelength}{ultraviolet}. $\qty{480}{nm}$~: visible,
bleu. $\qty{589}{nm}$~: visible, jaune-orange. $\qty{656}{nm}$~: visible,
rouge. $\qty{1064}{nm} > \qty{800}{nm}$~: \omterm{def:g10:light-spectra:wavelength}{infrarouge} --- ce qui rend un
laser de découpe doublement dangereux~: puissant et invisible.
\end{solution}

\begin{solution}{exo:g10:light-spectra:2}
\emph{1.} Une bande continue de couleurs, du rouge au violet~: le
\omterm{def:g10:light-spectra:white}{spectre} de la lumière solaire.

\emph{2.} L'extrémité violette --- le prisme dévie le plus les
\omterm{def:g10:light-spectra:wavelength}{courtes longueurs d'onde}.

\emph{3.} Une seule tache verte~: déviée, mais non étalée en bande. La
lumière laser est \omterm{def:g10:light-spectra:white}{monochromatique} --- une seule
\omterm{def:g10:light-spectra:wavelength}{longueur d'onde}, donc rien à trier.
\end{solution}

\begin{solution}{exo:g10:light-spectra:3}
(a) Continu (filament dense et chaud). (b) Émission de raies (gaz excité
à basse pression). (c) Continu (métal dense en fusion). (d) Émission de
raies (essentiellement la seule raie à $\qty{589}{nm}$). (e) Absorption~:
le \omterm{def:g10:light-spectra:continuous}{spectre continu} de la surface chaude, strié des raies sombres
de l'atmosphère plus froide du Soleil.
\end{solution}

\begin{solution}{exo:g10:light-spectra:4}
\emph{1.} Les gouttes de pluie~: chacune reçoit la lumière solaire
blanche, la décompose et la renvoie triée par couleur.

\emph{2.} Le Soleil fournit la \omterm{def:g10:light-spectra:white}{lumière blanche}~; la pluie fournit des
millions de petits prismes~; la géométrie ne fonctionne que pour la
lumière renvoyée par les gouttes vers l'observateur, c'est pourquoi le
Soleil doit être dans le dos de l'observateur et la pluie devant.

\emph{3.} Les gouttes n'ajoutent aucune lumière de leur propre --- elles
ne font que trier ce qu'elles reçoivent. Si la lumière solaire triée
affiche toutes les couleurs, alors toutes les couleurs étaient déjà
contenues dans la lumière solaire~: la lumière solaire est de la
\omterm{def:g10:light-spectra:white}{lumière blanche}.
\end{solution}

\begin{solution}{exo:g10:light-spectra:5}
\emph{1.} Rouge sombre $<$ orange vif $<$ presque blanc.

\emph{2.} Les deux~: la lueur devient \emph{plus brillante} (effet~1), et
sa couleur glisse du rouge vers le côté bleu (effet~2, le pic
$\lambda_{\max}$ se rétrécissant), de sorte qu'une part de plus en plus
grande de la bande visible est fortement éclairée --- une bande complète
se lit comme blanc.

\emph{3.} Un corps qui brille blanc émet fortement sur toute la bande
visible~; un corps qui brille rouge n'atteint que l'extrémité des
\omterm{def:g10:light-spectra:wavelength}{longues longueurs d'onde}. Par la loi de Wien le premier est le
plus chaud.
\end{solution}

\begin{solution}{exo:g10:light-spectra:6}
\emph{1.}
\[
T = \frac{\qty{2.90e-3}{m.K}}{\qty{5.00e-7}{m}} = \qty{5800}{K}.
\]

\emph{2.} $\theta = 5800 - 273 = \qty{5527}{\degreeCelsius} \approx
\qty{5500}{\degreeCelsius}$.

\emph{3.} $\qty{400}{nm} < \qty{500}{nm} < \qty{800}{nm}$~: le pic est
visible, dans le vert.
\end{solution}

\begin{solution}{exo:g10:light-spectra:7}
\emph{1.} Bougie~:
$\lambda_{\max} = \qty{2.90e-3}{m.K} / \qty{1800}{K} = \qty{1.61e-6}{m}
= \qty{1.6}{\micro m}$, dans l'\omterm{def:g10:light-spectra:wavelength}{infrarouge}. Plaque~:
$\lambda_{\max} = \qty{2.90e-3}{m.K} / \qty{500}{K} = \qty{5.8e-6}{m}
= \qty{5.8}{\micro m}$, loin dans l'\omterm{def:g10:light-spectra:wavelength}{infrarouge}.

\emph{2.} À $\qty{500}{K}$ pratiquement tout le rayonnement est
\omterm{def:g10:light-spectra:wavelength}{infrarouge}~: la peau l'absorbe et sent la chaleur, mais l'œil ne
reçoit rien --- chaud, pourtant noir.

\emph{3.} Un \omterm{def:g10:light-spectra:continuous}{spectre continu} est large~: même culminant à
$\qty{1.6}{\micro m}$, le \omterm{def:g10:light-spectra:white}{spectre} de la bougie garde une queue visible
du côté rouge-jaune --- faible, mais suffisante pour l'œil adapté à
l'obscurité.
\end{solution}

\begin{solution}{exo:g10:light-spectra:8}
Lampe~A~: les raies $434$, $486$, $\qty{656}{nm}$ correspondent à
l'hydrogène, et la faible raie violette complète le catalogue à
$\qty{410}{nm}$~: hydrogène.

Lampe~B~: la raie siège à $\qty{589}{nm}$, seule raie visible forte du
sodium. L'hélium est exclu non par la raie elle-même (l'hélium possède
bien une raie voisine à $\qty{588}{nm}$) mais par celles qui
\emph{manquent}~: l'hélium montrerait aussi $447$, $502$ et
$\qty{668}{nm}$, et aucune n'apparaît. Un élément s'identifie par son
catalogue complet~: la lampe~B contient du sodium.
\end{solution}

\begin{solution}{exo:g10:light-spectra:9}
\emph{1.} Un arc-en-ciel continu strié d'une fine raie sombre~: le
\omterm{def:g10:light-spectra:absorption}{spectre d'absorption} du sodium.

\emph{2.} À $\qty{589}{nm}$~: un gaz absorbe exactement les
\omterm{def:g10:light-spectra:wavelength}{longueurs d'onde} qu'il est capable d'émettre
(\cref{prop:g10:light-spectra:kirchhoff}), et la raie visible forte du
sodium est à $\qty{589}{nm}$.

\emph{3.} Le \omterm{def:g10:light-spectra:emission}{spectre d'émission} du sodium~: une seule raie brillante
jaune-orange à $\qty{589}{nm}$ sur fond noir --- la même
\omterm{def:g10:light-spectra:wavelength}{longueur d'onde}, les deux spectres échangeant brillant et sombre.
\end{solution}

\begin{solution}{exo:g10:light-spectra:10}
\emph{1.} $410$, $434$, $486$, $\qty{656}{nm}$~: le catalogue complet de
l'hydrogène. $517$ et $\qty{518}{nm}$~: la paire du magnésium. Les deux
éléments sont présents.

\emph{2.} Dans l'atmosphère du Soleil, le gaz plus froid et plus mince
au-dessus de la surface lumineuse. La surface envoie un
\omterm{def:g10:light-spectra:continuous}{spectre continu}~; l'atmosphère retire ses propres
\omterm{def:g10:light-spectra:wavelength}{longueurs d'onde} de la lumière qui la traverse, donc les raies
apparaissent comme de la lumière manquante --- sombres contre le fond
brillant.
\end{solution}

\begin{solution}{exo:g10:light-spectra:11}
\emph{1.} $T = 37 + 273 = \qty{310}{K}$, donc
\[
\lambda_{\max} = \frac{\qty{2.90e-3}{m.K}}{\qty{310}{K}}
= \qty{9.4e-6}{m} = \qty{9.4}{\micro m}.
\]

\emph{2.} Loin dans l'\omterm{def:g10:light-spectra:wavelength}{infrarouge} --- plus de dix fois la
\omterm{def:g10:light-spectra:wavelength}{longueur d'onde} de la lumière rouge.

\emph{3.} À $\qty{310}{K}$ la part visible du rayonnement est
totalement négligeable~: aucune lueur visible, si sombre que soit la
pièce. Une caméra thermique porte un capteur accordé sur le rayonnement
réellement émis, autour de $\qty{9}{\micro m}$, et voit les corps
chauds briller contre un environnement plus froid.
\end{solution}

\begin{solution}{exo:g10:light-spectra:12}
\emph{1.} $\lambda_{\max} = \qty{2.90e-3}{m.K} / \qty{2700}{K}
= \qty{1.07e-6}{m} \approx \qty{1.1}{\micro m}$~: proche
\omterm{def:g10:light-spectra:wavelength}{infrarouge}.

\emph{2.} Le \omterm{def:g10:light-spectra:white}{spectre} culmine au-delà du visible~: la plus grande partie
de la puissance électrique revient en \omterm{def:g10:light-spectra:wavelength}{infrarouge} invisible ---
c'est-à-dire en chaleur --- et seule une faible queue visible éclaire la
pièce. Comme lampe l'ampoule gaspille la plupart de sa puissance~; comme
radiateur elle est presque parfaite.

\emph{3.} $T = \qty{2.90e-3}{m.K} / \qty{5.50e-7}{m} \approx
\qty{5300}{K}$ --- bien au-dessus du point de fusion du tungstène de
$\qty{3700}{K}$. Aucun filament solide ne peut être poussé aux
températures de la lumière du jour, donc les lampes de type «~jour~»
ont dû abandonner complètement l'incandescence (tubes fluorescents,
LED), produisant de la lumière visible sans chauffer un corps à
$\qty{5300}{K}$.
\end{solution}

\begin{solution}{exo:g10:light-spectra:13}
\emph{1.} $T = \qty{2.90e-3}{m.K} / \qty{4.20e-7}{m} \approx
\qty{6900}{K}$~: plus chaude que le Soleil à $\qty{5800}{K}$.

\emph{2.} Hydrogène (l'ensemble complet $410$, $434$, $486$,
$\qty{656}{nm}$ est présent) et calcium (les deux raies, $393$ et
$\qty{397}{nm}$).

\emph{3.} Non. La raie du sodium à $\qty{589}{nm}$ est absente, alors
qu'un sodium abondant dans l'atmosphère l'imprimerait nécessairement sur
le fond continu. Si le sodium est présent, ce n'est qu'en traces trop
faibles pour être détectées.
\end{solution}

\begin{solution}{exo:g10:light-spectra:14}
\emph{1.} L'hydrogène explique $410$, $434$, $486$ et $\qty{656}{nm}$~;
l'hélium explique $447$, $502$, $588$ et $\qty{668}{nm}$. Chaque raie
observée est comptabilisée~: le mélange est hydrogène et hélium.

\emph{2.} L'hélium est une explication \emph{suffisante} parce que ses
trois autres raies ($447$, $502$, $\qty{668}{nm}$) sont toutes
présentes --- la raie à $\qty{588}{nm}$ est alors attendue de l'hélium
seul, et rien dans le \omterm{def:g10:light-spectra:white}{spectre} n'\emph{exige} le sodium. Pour trancher
il faut de la résolution, non de l'argument~: un \omterm{def:g10:light-spectra:white}{spectroscope} plus
fin sépare la raie unique de l'hélium à $\qty{587.6}{nm}$ de celle du
sodium, qui est en fait une paire rapprochée à $589.0$ et
$\qty{589.6}{nm}$. Voir une seule raie là-bas absout le sodium~; en
voir trois le confond.
\end{solution}

\begin{solution}{exo:g10:light-spectra:15}
\emph{1.} Bételgeuse~:
$T = \qty{2.90e-3}{m.K} / \qty{8.5e-7}{m} \approx \qty{3400}{K}$. Rigel~:
$T = \qty{2.90e-3}{m.K} / \qty{2.4e-7}{m} \approx \qty{12000}{K}$.

\emph{2.} Comme $\lambda_{\max} \times T$ est la même constante pour les
deux,
$\dfrac{T_{\text{Rigel}}}{T_{\text{Betelgeuse}}}
= \dfrac{\qty{850}{nm}}{\qty{240}{nm}} \approx 3.5$.

\emph{3.} Bételgeuse culmine dans l'\omterm{def:g10:light-spectra:wavelength}{infrarouge}~: dans la bande
visible son \omterm{def:g10:light-spectra:white}{spectre} est le plus fort à l'extrémité rouge, donc
l'étoile paraît rouge. Rigel culmine dans l'\omterm{def:g10:light-spectra:wavelength}{ultraviolet}~: sa bande
visible est la plus forte à l'extrémité bleue, d'où le bleu-blanc.

\emph{4.} Un \omterm{def:g10:light-spectra:continuous}{spectre continu} couvre un énorme intervalle de
\omterm{def:g10:light-spectra:wavelength}{longueurs d'onde}. Même avec son pic hors de la bande visible,
une étoile verse abondamment de lumière \emph{dans} cette bande --- le
pic décide de la teinte, non de la visibilité.
\end{solution}

\begin{solution}{pb:g10:light-spectra:1}
\textbf{1.} D'environ $\qty{400}{nm}$ (extrémité violette) à environ
$\qty{800}{nm}$ (extrémité rouge)~; $\qty{550}{nm}$ se trouve près du
milieu, dans le vert.

\textbf{2.} $\qty{410}{nm}$~: violet~; $\qty{434}{nm}$~: violet-bleu~;
$\qty{486}{nm}$~: bleu-vert~; $\qty{656}{nm}$~: rouge.

\textbf{3.} Une lueur \omterm{def:g10:light-spectra:white}{monochromatique} jaune-orange. La couleur d'une
voiture est le mélange des \omterm{def:g10:light-spectra:wavelength}{longueurs d'onde} que sa peinture
réfléchit~; sous une lampe n'offrant \emph{que} $\qty{589}{nm}$, chaque
peinture ne peut que réfléchir plus ou moins de cette seule
\omterm{def:g10:light-spectra:wavelength}{longueur d'onde} --- toutes les voitures apparaissent en nuances
d'orange et de gris.

\textbf{4.} Chaque élément émet un catalogue fixe de raies, et aucun
élément ne partage le même catalogue
(\cref{prop:g10:light-spectra:fingerprint})~: un ensemble complet de
raies identifie donc l'élément comme une empreinte digitale. Mais une
correspondance \emph{partielle} prouve peu --- des éléments distincts
peuvent posséder des raies isolées quasi coïncidentes (hélium à
$\qty{588}{nm}$, sodium à $\qty{589}{nm}$) --- donc la prétention de
présence exige toutes les raies fortes du catalogue.

\textbf{5.} (a) Un arc-en-ciel continu avec une raie sombre à
$\qty{589}{nm}$~: le \omterm{def:g10:light-spectra:absorption}{spectre d'absorption}. (b) Une seule raie
brillante à $\qty{589}{nm}$ sur fond noir~: le \omterm{def:g10:light-spectra:emission}{spectre d'émission}
--- même \omterm{def:g10:light-spectra:wavelength}{longueur d'onde}, brillant et sombre échangés
(\cref{prop:g10:light-spectra:kirchhoff}).

\textbf{6.} Le \omterm{def:g10:light-spectra:white}{spectre} s'éclaircit à toute \omterm{def:g10:light-spectra:wavelength}{longueur d'onde}, et son
pic $\lambda_{\max}$ se déplace vers les \omterm{def:g10:light-spectra:wavelength}{courtes longueurs d'onde},
suivant $\lambda_{\max} \times T = \qty{2.90e-3}{m.K}$
(\cref{prop:g10:light-spectra:wien}).

\textbf{7.} $\lambda_{\max} = \qty{2.90e-3}{m.K} / \qty{3400}{K}
= \qty{8.5e-7}{m} = \qty{850}{nm}$~: proche \omterm{def:g10:light-spectra:wavelength}{infrarouge}, juste
après le bord rouge --- c'est pourquoi la lumière halogène est plus
chaude de ton que la lumière du jour.

\textbf{8.} $T = \qty{2.90e-3}{m.K} / \qty{5.00e-7}{m} = \qty{5800}{K}$.

\textbf{9.} Le pic est doux, non un pic aigu~: le Soleil émet toute la
bande visible avec une force comparable. Une bande complète de couleurs
mélangées se lit comme blanc --- le Soleil paraît blanc (jaunâtre à
travers notre atmosphère), jamais vert.

\textbf{10.} Tison froid~: pic profond dans l'\omterm{def:g10:light-spectra:wavelength}{infrarouge}, aucune
queue visible --- le rayonnement se sent, rien ne se voit. Plus chaud~:
le \omterm{def:g10:light-spectra:white}{spectre} qui s'éclaircit pousse une première queue visible au-delà
de $\qty{800}{nm}$ --- l'extrémité rouge s'allume seule~: rouge sombre.
Plus chaud encore~: toute la bande visible se remplit tandis que tout
s'éclaircit~: orange, puis vers le blanc.

\textbf{11.} Le fond continu vient de la surface dense et chaude qui
brille~; les raies sombres sont imprimées par l'atmosphère plus froide
et plus mince au-dessus, qui absorbe ses propres
\omterm{def:g10:light-spectra:wavelength}{longueurs d'onde} dans la lumière qui la traverse.

\textbf{12.} Hydrogène ($410$, $434$, $486$, $\qty{656}{nm}$~: complet),
sodium ($\qty{589}{nm}$) et calcium ($393$ et $\qty{397}{nm}$~: les
deux) --- les trois sont présents dans l'atmosphère du Soleil.

\textbf{13.} Que l'atmosphère est dominée par l'hydrogène. Les mesures
modernes s'accordent et quantifient~: le Soleil est environ aux trois
quarts de l'hydrogène en masse, le reste étant surtout de l'hélium.

\textbf{14.} La raie ne correspondait à aucun élément connu~; comme le
catalogue de chaque élément est fixe et unique, une raie n'appartenant
à aucun catalogue doit appartenir à un élément jamais encore vu en
laboratoire. Un nouvel élément fut annoncé et nommé hélium, d'après
\emph{helios}, le Soleil. Il fut isolé sur Terre en 1895 --- la
prédiction confirmée, vingt-sept ans plus tard.

\textbf{15.} Vue contre le fond continu éclatant, l'atmosphère
\emph{retire} de la lumière à ses propres \omterm{def:g10:light-spectra:wavelength}{longueurs d'onde}~: raies
sombres. Pendant l'éclipse la Lune bloque le fond~; la mince
atmosphère se dresse alors seule contre un ciel sombre et sa lumière
réémise est tout ce qu'il y a à voir~: les mêmes
\omterm{def:g10:light-spectra:wavelength}{longueurs d'onde}, maintenant brillantes. Un gaz, un catalogue
--- le signe dépend de ce qui se tient derrière.

\textbf{16.} $T = \qty{2.90e-3}{m.K} / \qty{2.90e-7}{m}
= \qty{10000}{K}$.

\textbf{17.} Bleu-blanc~: dans la bande visible le \omterm{def:g10:light-spectra:white}{spectre} est le
plus fort à l'extrémité bleue. Et pas d'invisibilité~: le
\omterm{def:g10:light-spectra:continuous}{spectre continu} inonde toute la bande visible sur le chemin du
pic \omterm{def:g10:light-spectra:wavelength}{ultraviolet} --- le pic fixe la teinte, non la brillance.

\textbf{18.} Les raies sombres $410$, $434$, $486$, $\qty{656}{nm}$ sont
l'empreinte complète de l'hydrogène~: l'hydrogène domine l'atmosphère
de Véga.

\textbf{19.} $T_{\text{Vega}} / T_{\text{Sun}} = 10000 / 5800 \approx
1.7$~: Véga est plus chaude, donc plus bleue (pic à $\qty{290}{nm}$
contre $\qty{500}{nm}$) et, par la première partie
du~\cref{prop:g10:light-spectra:wien}, plus brillante par
\omterm{def:g10:orders-of-magnitude:unit}{unité} de surface rayonnante à toute
\omterm{def:g10:light-spectra:wavelength}{longueur d'onde}.

\textbf{20.} Fiche d'identité de Véga --- température de surface~:
environ $\qty{10000}{K}$~; élément dominant de l'atmosphère~:
hydrogène~; couleur apparente~: bleu-blanc. Le tout lu dans un seul
faisceau de lumière, avec un prisme, un catalogue de raies de
laboratoire et une division. Le «~jamais~» de Comte fut émis en 1835
et expira en 1859~: il dura vingt-quatre ans.
\end{solution}
